\section{Recursive colors and finite averages}\label{sec:proof}

We prove Theorem~\ref{thm:main} using the map supplied by
Lemma~\ref{lem:displacement}.  The argument uses a single finite average
over the group.  We first record the dyadic facts that make its
correlations comparable.

\subsection{Dyadic partitions and affine transport}\label{sec:dyadic}

We use the standard dyadic-partition model of $F$; see
\cite[Section~2, Lemma~2.2]{CannonFloydParry1996}. We record the
restriction and pair-transport facts needed for the finite averages.

A \emph{basic dyadic interval} is an interval
\[
  [k2^{-r},(k+1)2^{-r}]\subseteq[0,1],
  \qquad r\in\N\cup\{0\},\quad 0\leq k<2^r.
\]
A \emph{basic partition} is a finite partition of $[0,1]$ into such
intervals; cells may share endpoints.  Its mesh is the greatest length
of a cell.  A partition $T$ \emph{respects} an interval $I$ if $I$ is a
union of cells of $T$.

For a basic interval $I$, let $s_I:[0,1]\to I$ be the increasing affine
map, and write $I\cdot J=s_I(J)$.  If $T$ respects $I$, define its
normalized restriction by
\[
  T_I=\{s_I^{-1}(K):K\in T,\ K\subseteq I\}.
\]
Both $I\cdot J$ and the cells of $T_I$ are basic dyadic intervals.
Indeed, if $I$ has length $2^{-r}$, a basic cell contained in $I$ has
length $2^{-s}$ with $s\geq r$; applying $s_I^{-1}$ gives a cell of
length $2^{-(s-r)}$ whose left endpoint is a multiple of that length.
The affine charts satisfy $s_{I\cdot J}=s_I\circ s_J$.  Consequently,
whenever $T$ respects both $I$ and $I\cdot J$, the partition $T_I$
respects $J$ and
\begin{equation}\label{dyad:restriction}
  (T_I)_J=T_{I\cdot J}.
\end{equation}

Let $T^{(n)}$ be the uniform partition of $[0,1]$ into intervals of
length $2^{-n}$.  For $g\in F$, the notation $gT^{(n)}$ denotes the
partition formed by the images of these cells.

\begin{lemma}\label{dyad:eventual}
For each $g\in F$, the partitions $gT^{(n)}$ are basic for all
sufficiently large $n$, and their meshes tend to zero.  In particular,
given a finite set $K\subset F$ and a finite family $\mathcal J$ of
basic intervals, there is an $n_0$ such that $gT^{(n)}$ is basic and
respects every member of $\mathcal J$ whenever $g\in K$ and $n\geq n_0$.
\end{lemma}

\begin{proof}
The image under $g$ of each breakpoint is dyadic: starting from
$g(0)=0$, sum the increments on the preceding linear pieces, each a
power of two times a dyadic length.  Thus every linear piece has the
form $x\mapsto 2^q x+b$ with $q\in\mathbb Z$ and $b$ dyadic.  Choose
$n$ large enough to resolve all breakpoints, to have $n\geq q$ on
every piece, and to make each corresponding $b$ an integer multiple
of $2^{q-n}$.  The image of each uniform cell is then a basic interval
of length $2^{q-n}$.  There are finitely many slopes, so the mesh tends
to zero.

Two basic intervals with intersecting interiors are nested.  Hence
a basic partition of mesh smaller than the length of a basic
interval $I$ respects $I$: every cell whose interior meets that of
$I$ must be contained in $I$.  Apply this observation to the
finitely many intervals in $\mathcal J$, and then take the maximum
of the finitely many thresholds for $g\in K$.
\end{proof}

Write $I<J$ when $\max I<\min J$; in particular this notation requires
a positive gap.  An interval is \emph{internal} if both its endpoints
lie in $(0,1)$.

\begin{lemma}\label{dyad:transport}
If $I<J$ and $I'<J'$ are pairs of internal basic intervals, there is
an $h\in F$ whose restrictions carry $I$ affinely onto $I'$ and $J$
affinely onto $J'$.
\end{lemma}

\begin{proof}
The two selected intervals leave three complementary gaps, all of
positive length and with dyadic endpoints.  Partition each gap into
basic cells, for example by a sufficiently fine uniform dyadic
grid.  Do this for both pairs.  Within each corresponding pair of
gaps, equalize the numbers of cells by successively bisecting cells
on the side with fewer cells.  Each bisection increases the number
by one.  The source and target partitions now have the same number
of cells, with the two selected intervals in matching positions.
Map corresponding cells increasingly and affinely.  All
breakpoints are dyadic and all slopes are ratios of basic lengths,
hence powers of two.  This defines the required element of $F$.
\end{proof}

The relevant covariance is exact.  If $h$ carries $I$ affinely onto
$I'$, then $h\circ s_I=s_{I'}$.  Thus, if $gT^{(n)}$ and
$(hg)T^{(n)}$ are basic and respect $I$ and $I'$, respectively,
\begin{equation}\label{dyad:covariance}
  \bigl((hg)T^{(n)}\bigr)_{I'}=(gT^{(n)})_I.
\end{equation}
Here, as throughout, $hg=h\circ g$.  To verify the equality, match
each source cell $C$ with its image $h(C)$ and apply
$s_{I'}^{-1}\circ h=s_I^{-1}$ on $I$.

\subsection{The recursive coloring}\label{sec:coloring}

Fix a real Hilbert space $H$, its closed unit ball $B$, and a
Lipschitz map $f:B\to B$ as in Lemma~\ref{lem:displacement}.  Fix
$L>0$ and $\delta>0$ such that
\begin{equation}\label{color:map}
  \norm{f(x)-f(y)}\leq L\norm{x-y},\qquad
  \norm{f(x)-x}\geq\delta
  \quad(x,y\in B).
\end{equation}
Choose an integer $D\geq2$ so that
\begin{equation}\label{color:parameters}
  D>\frac{4L^2}{\delta^2}.
\end{equation}
Choose internal basic intervals $I_1<\cdots<I_D$.  Such a family
exists for every finite $D$: take $2^r>2D$ and use
$I_j=[(2j-1)2^{-r},2j2^{-r}]$.

Define a color $p(T)\in B$ for every basic partition $T$ by induction
on its number of cells:
\begin{equation}\label{color:recursion}
  p(T)=
  \begin{cases}
  \displaystyle f\left(\frac1D\sum_{j=1}^D p(T_{I_j})\right),
       &\text{if $T$ respects every $I_j$},\\[6pt]
  0,   &\text{otherwise}.
  \end{cases}
\end{equation}
This induction is well founded.  In the first case each $T_{I_j}$
has strictly fewer cells than $T$, since $I_j$ is internal and $T$
has cells outside it.  The average of the previously defined
colors belongs to the convex ball $B$, so the argument of $f$ is
always in its domain, and $p(T)\in B$.

For the displayed choice $|I_j|=2^{-r}$, the one-cell partition
$T^{(0)}$ has color zero, while
\[
  p(T^{(r)})=f(0),\qquad
  p(T^{(2r)})=f(f(0)).
\]
Indeed, normalizing a restriction to any $I_j$ changes these uniform
partition levels from $r$ to $0$ and from $2r$ to $r$, respectively.

We use the finite family of parent intervals and their descendants
\[
  \mathcal I=\{I_i:1\leq i\leq D\}
       \cup\{I_i\cdot I_j:1\leq i,j\leq D\}.
\]
All these intervals are internal.  Call $(n,g)$ \emph{admissible}
if $gT^{(n)}$ is a basic partition respecting every member of
$\mathcal I$.  For every $n$, every $g\in F$, and every
$I\in\mathcal I$, define
\begin{equation}\label{color:global}
  X_I^{(n)}(g)=
  \begin{cases}
    p\bigl((gT^{(n)})_I\bigr),&\text{if $(n,g)$ is admissible},\\
    0,&\text{otherwise}.
  \end{cases}
\end{equation}
In particular these are globally defined $B$-valued functions.
Their scalar correlations
\[
  \varphi_{I,J}^{(n)}(g)
     =\ip{X_I^{(n)}(g)}{X_J^{(n)}(g)}
\]
are defined on all of $F$ and have absolute value at most one.
By Lemma~\ref{dyad:eventual}, every finite set of group elements is
admissible at one common sufficiently large level $n$.

\subsection{A uniform boundary bound}\label{sec:boundary}

For every
ordered pair $I<J$ in $\mathcal I$, choose, using
Lemma~\ref{dyad:transport}, an element $h_{I,J}\in F$ carrying
$I,J$ affinely onto $I_1,I_2$, respectively.  Let $S$ be the finite
set of these elements.  All choices so far, in particular $S$, are
independent of the finite set to be tested.

For any nonempty finite \(A\subset F\), any \(h\in F\), and any
\(\varphi:F\to[-1,1]\), cancellation over \(hA\cap A\) gives
\begin{equation}\label{eq:finite-translation}
 \left|\frac1{|A|}\sum_{g\in A}\varphi(hg)
       -\frac1{|A|}\sum_{g\in A}\varphi(g)\right|
 \le\frac{|hA\mathbin{\triangle}A|}{|A|}.
\end{equation}
This is the finite-average comparison we will apply to the scalar
correlations.

\begin{proposition}\label{avg:boundary}
For every nonempty finite set $A\subset F$,
\begin{equation}\label{avg:boundary-bound}
  \max_{h\in S}\frac{|hA\mathbin{\triangle}A|}{|A|}
  \geq \frac{\delta^2/L^2-4/D}{4(1-1/D)}>0.
\end{equation}
\end{proposition}

\begin{proof}
Fix such a set $A$ and put
\begin{equation}\label{avg:defect}
  \eta=\max_{h\in S}\frac{|hA\mathbin{\triangle}A|}{|A|}.
\end{equation}
Choose a single $n$ for which $(n,g)$ is admissible for all
\[
  g\in A\cup SA,
  \qquad SA=\bigcup_{h\in S}hA.
\]
This choice is possible because that union is finite; the level $n$
may depend on $A$.  We now fix this $n$, omit its superscript, and write
\[
  \mathbb E_A\psi=\frac1{|A|}\sum_{g\in A}\psi(g),
  \qquad
  \alpha=\mathbb E_A\varphi_{I_1,I_2}.
\]
\paragraph{Comparing correlations.}
For $I<J$ in $\mathcal I$ and $g\in A$, admissibility of both $g$
and $h_{I,J}g$, together with \eqref{dyad:covariance}, gives
\[
  X_I(g)=X_{I_1}(h_{I,J}g),\qquad
  X_J(g)=X_{I_2}(h_{I,J}g).
\]
It follows that $\varphi_{I,J}(g)
=\varphi_{I_1,I_2}(h_{I,J}g)$.  Apply
\eqref{eq:finite-translation} to the globally bounded function
$\varphi_{I_1,I_2}$ and use \eqref{avg:defect}.  We obtain
\begin{equation}\label{avg:correlation}
  \left|\mathbb E_A\ip{X_I}{X_J}-\alpha\right|
       \leq\eta
  \qquad\text{whenever $I,J\in\mathcal I$ are strictly separated}.
\end{equation}
The reversed order follows from symmetry of the real inner
product.  No condition is imposed on nested pairs.  Also, a
transport $h_{I,J}$ need not preserve the other intervals in
$\mathcal I$: admissibility was required separately at its source
and target group elements.

\paragraph{The variance estimate.}

For $g\in A$, put
\[
  m(g)=\frac1D\sum_{k=1}^D X_{I_k}(g),\qquad
  z_i(g)=\frac1D\sum_{j=1}^D X_{I_i\cdot I_j}(g)
       \quad(1\leq i\leq D).
\]
These vectors lie in $B$.  Admissibility ensures that the partition
$(gT^{(n)})_{I_i}$ respects each $I_j$.  Applying the recursion and
\eqref{dyad:restriction} therefore yields the exact identities
\begin{equation}\label{avg:recursive-identity}
  X_{I_i}(g)=f(z_i(g))\qquad(g\in A, 1\leq i\leq D).
\end{equation}
Consequently, for every $g\in A$, the displacement bound, the
convexity of the squared norm, and the Lipschitz bound give
\begin{align}
  \delta^2
  &\leq\norm{m(g)-f(m(g))}^2\notag\\
  &=\left\|\frac1D\sum_{i=1}^D
                   \bigl(f(z_i(g))-f(m(g))\bigr)\right\|^2\notag\\
  &\leq\frac1D\sum_{i=1}^D
                   \norm{f(z_i(g))-f(m(g))}^2\notag\\
  &\leq\frac{L^2}{D}\sum_{i=1}^D\norm{z_i(g)-m(g)}^2.
       \label{avg:displacement}
\end{align}

We bound the finite average of each squared distance in the last
line.  In the expansion of $\norm{m}^2$, the $D(D-1)$
off-diagonal terms pair strictly separated parents, so
\eqref{avg:correlation} applies.  Each of the $D$ diagonal terms
is at most one.  The same statements hold for the siblings in
$\norm{z_i}^2$.  Hence
\begin{align}
  \mathbb E_A\norm{m}^2
       &\leq\left(1-\frac1D\right)(\alpha+\eta)+\frac1D,
       \label{avg:parent-bound}\\
  \mathbb E_A\norm{z_i}^2
       &\leq\left(1-\frac1D\right)(\alpha+\eta)+\frac1D.
       \label{avg:child-bound}
\end{align}
In the mixed term
\[
  \ip{z_i}{m}
    =\frac1{D^2}\sum_{j=1}^D\sum_{k=1}^D
                  \ip{X_{I_i\cdot I_j}}{X_{I_k}},
\]
the $D(D-1)$ terms with $k\ne i$ pair strictly separated
intervals.  Each has average at least $\alpha-\eta$.
The remaining $D$ terms pair a descendant with its own parent;
each is at least $-1$ by the unit-ball bound.  Figure~\ref{fig:mixed-pairs}
shows this division of the mixed pairs.  Thus
\begin{equation}\label{avg:mixed-bound}
  \mathbb E_A\ip{z_i}{m}
    \geq\left(1-\frac1D\right)(\alpha-\eta)-\frac1D.
\end{equation}
\begin{figure}[tbp]
\centering
\begin{tikzpicture}[x=1.05cm,y=0.62cm,font=\small]
  \foreach \x in {0,1,2,3,4} {
    \foreach \y in {0,1,2,3} {
      \fill[parentblue!13] (\x,\y) rectangle +(1,1);
    }
  }
  \fill[childorange!24] (2,0) rectangle (3,4);
  \draw[gray!65,thin] (0,0) grid[xstep=1,ystep=1] (5,4);
  \draw[childorange!65!black,semithick] (2,0) rectangle (3,4);
  \node at (0.5,4.45) {$I_1$};
  \node at (1.5,4.45) {$\cdots$};
  \node[childorange] at (2.5,4.45) {$I_i$};
  \node at (3.5,4.45) {$\cdots$};
  \node at (4.5,4.45) {$I_D$};
  \node[anchor=east] at (-0.2,3.5) {$I_i\cdot I_1$};
  \node[anchor=east] at (-0.2,2.5) {$\vdots$};
  \node[anchor=east] at (-0.2,1.5) {$I_i\cdot I_j$};
  \node[anchor=east] at (-0.2,0.5) {$I_i\cdot I_D$};
  \node[anchor=west,align=left,text width=5.2cm] at (5.45,3.15)
    {\textcolor{parentblue}{Separated pairs}\\
     \(D(D-1)\) terms, each with\\
     mean in \([\alpha-\eta,\alpha+\eta]\)};
  \node[anchor=west,align=left,text width=5.2cm] at (5.45,0.95)
    {\textcolor{childorange}{Nested pairs}\\
     \(D\) terms, each at least \(-1\)};
\end{tikzpicture}
\caption{For fixed \(i\), the \((j,k)\) entry is
\(\ip{X_{I_i\cdot I_j}}{X_{I_k}}\). Rows index the children of
\(I_i\) and columns index all parents; the grid is schematic.
Only column \(i\) contains nested pairs, a fraction \(1/D\) of the
\(D^2\) terms.}
\label{fig:mixed-pairs}
\end{figure}

Expanding the squared distance and using
\eqref{avg:parent-bound}--\eqref{avg:mixed-bound}, the coefficients
of $\alpha$ cancel and give
\begin{equation}\label{avg:variance}
  \mathbb E_A\norm{z_i-m}^2
    \leq\frac4D+4\left(1-\frac1D\right)\eta
    \qquad(1\leq i\leq D).
\end{equation}
This estimate does not require any sign assumption on $\alpha$.

Average \eqref{avg:displacement} over $A$ and apply
\eqref{avg:variance}:
\begin{align*}
  \delta^2
  &\leq\frac{L^2}{D}\sum_{i=1}^D
                 \mathbb E_A\norm{z_i-m}^2\\
  &\leq L^2\left(\frac4D+4\left(1-\frac1D\right)\eta\right).
\end{align*}
Rearranging proves \eqref{avg:boundary-bound}; its right-hand side is
positive by \eqref{color:parameters} and does not depend on $A$.
\end{proof}

\begin{proof}[Proof of Theorem~\ref{thm:main}]
The fixed finite set $S$ and the uniform positive lower bound in
Proposition~\ref{avg:boundary} violate the F\o lner criterion
\eqref{eq:folner}.  Hence $F$ is not amenable.
\end{proof}
